% Part: first-order-logic
% Chapter: lindstrom
% Section: abstract-logics

\documentclass[../../../include/open-logic-section]{subfiles}

\begin{document}

\olfileid{mod}{lin}{alg}
\olsection{Logika Abstrak}

\begin{defn}
Sawijining \emph{logika abstrak} yaiku pasangan $\tuple{L, \models_L}$,
kanthi $L$ iku fungsi kang masangake marang saben !!{language}~$\Lang{L}$
sawijining himpunan $L(\Lang{L})$ saka !!{sentence}s, lan $\models_L$
iku relasi antarane !!{structure}s kanggo !!{language}~$\Lang{L}$
lan !!{element}s saka $L(\Lang{L})$. Mligine, $\tuple{F, \models}$
iku logika orde kapisan lumrah, yaiku $F$ fungsi kang masangake marang
!!{language}~$\Lang{L}$ himpunan !!{sentence}s orde kapisan kang
kawangun saka konstanta ing $\Lang{L}$, lan $\models$ iku relasi
nyukupi ing logika orde kapisan.
\end{defn}

Gatekna yen kita isih nggunakake pangerten !!{structure} kang padha
kanggo !!{language} tartamtu kaya ing logika orde kapisan, nanging kita
ora nganggep yen !!{sentence}s kawangun saka simbol-simbol dhasar ing
$\Lang{L}$ kanthi cara lumrah, utawa yen relasi $\models_L$
ditetepake kanthi rekursif kaya ing logika orde kapisan. Mula, umpamane,
definisi kang sipate umum iki uga nyakup !!{sentence}s ing
$\tuple{L,\models_L}$ kang ngemot konjungsi utawa disjungsi tanpa wates
dawane, kuantor saliyane $\lexists$ lan $\lforall$ (umpamane,
“ana tanpa wates akehe~$x$ saengga \dots”), utawa mbokmenawa
runtutan kuantor tanpa wates dawane. Kanggo negesake yen
“!!{sentence}s” ing $L(\Lang{L})$ ora kudu awujud !!{sentence}s
lumrah ing logika orde kapisan, ing bab iki kita nggunakake
!!{variable}s $!E$, $!F$,~\dots kanggo nglambangake ukara-ukara iku,
lan nyisihake $!A$, $!B$,~\dots kanggo !!{formula}s lumrah ing
logika orde kapisan.

\begin{defn}
Upamana $\Mod(L){!E}$ nglambangake kelas
$\Setabs{\Struct{M}}{\Struct{M} \models_L !E}$. Yen !!{language}
kudu dicethakake, kita nulis $\Mod[L](L){!E}$. Rong !!{structure}s
$\Struct{M}$ lan $\Struct{N}$ kanggo $\Lang{L}$ iku
\emph{ekuivalen elementer ing} $\tuple{L, \models_L}$, ditulis
$\Struct{M} \elemequiv[L] \Struct{N}$, yen !!{sentence}s kang
padha saka $L(\Lang{L})$ bener ing kalorone.
\end{defn}

\begin{defn}
Sawijining logika abstrak $\tuple{L,\models_L}$ kanggo
!!{language} $\Lang{L}$ diarani \emph{normal} yen nyukupi
sipat-sipat ing ngisor iki:
\begin{enumerate}
\item (\emph{Monotonisitas-$L$}) Kanggo !!{language}s $\Lang{L}$
  lan $\Lang{L'}$, yen $\Lang{L} \subseteq \Lang{L'}$, mula
  $L(\Lang{L}) \subseteq L(\Lang{L'})$.
\item (\emph{Sipat Ekspansi}) Kanggo saben $!E \in L(\Lang{L})$
  ana subhimpunan \emph{winates} $\Lang{L'}$ saka $\Lang{L}$
  saengga relasi $\Struct{M} \models_L !E$ mung gumantung marang
  redukt $\Struct{M}$ menyang $\Lang{L'}$; yaiku, yen
  $\Struct{M}$ lan $\Struct{N}$ nduweni redukt kang padha menyang
  $\Lang{L'}$, mula $\Struct{M} \models_L !E$ yen lan mung yen
  $\Struct{N} \models_L !E$.
\item (\emph{Sipat Isomorfisme}) Yen $\Struct{M} \models_L !E$
  lan $\Struct{M} \simeq \Struct{N}$, mula $\Struct{N} \models_L
  !E$ uga.
\item (\emph{Sipat Pangowahan Jeneng}) Relasi $\models_L$ tetep
  lumaku nalika simbol-simbol diganti jenenge: yen !!{language}
  $\Lang{L}'$ dipikolehi saka $\Lang{L}$ kanthi ngganti saben
  simbol $P$ nganggo simbol $P'$ kang aritase padha lan saben
  konstanta $c$ nganggo konstanta $c'$ kang beda, mula kanggo
  saben !!{structure}~$\Struct{M}$ lan !!{sentence}~$!E$,
  $\Struct{M} \models_L !E$ yen lan mung yen
  $\Struct{M}' \models_L !E'$, kanthi $\Struct{M}'$ iku
  !!{structure} kanggo $\Lang{L}'$ kang cocog karo
  $\Lang{L}$ lan $!E' \in L(\Lang{L}')$.
\item (\emph{Sipat Boolean}) Logika abstrak $\tuple{L,
  \models_L}$ katutup tumrap konektif Boolean ing teges iki:
  kanggo saben $!E \in L(\Lang{L})$ ana $!F \in
  L(\Lang{L})$ saengga $\Struct{M} \models_L !F$ yen lan mung
  yen $\Struct{M} \not\models_L !E$, lan kanggo saben $!E$
  lan $!F$ ana $!G$ saengga $\Mod(L){!G} = \Mod(L){!E}
  \cap \Mod(L){!F}$. Semono uga kanggo !!{formula}s atomik
  lan konektif liyane.
\item (\emph{Sipat Kuantor}) Kanggo saben konstanta $c$ ing
  $\Lang{L}$ lan $!E \in L(\Lang{L})$, ana $!F \in
  L(\Lang{L})$ saengga
  \[
  \Mod[L'](L){!F} = \Setabs{\Struct{M}}{\Expan{M}{a} \in
  \Mod[L](L){!E} \text{ kanggo sawatara } a \in \Domain{M}},
  \]
  kanthi $\Lang{L'} = \Lang{L} \setminus \{c\}$ lan
  $\Expan{M}{a}$ iku ekspansi $\Struct{M}$ menyang $\Lang{L}$
  kang masangake $a$ marang~$c$.
\item (\emph{Sipat Relativisasi}) Yen diwenehi !!a{sentence}
  $!E \in L(\Lang{L})$ lan simbol-simbol $R$, $c_1$, \dots,
  $c_n$ kang ora ana ing $\Lang{L}$, ana !!a{sentence}
  $!F \in L(\Lang{L} \cup \{R,c_1,\ldots,c_n\})$
  kang diarani \emph{relativisasi} $!E$ marang
  $\Atom{R}{x, c_1, \dots c_n}$, saengga kanggo saben
  !!{structure}~$\Struct{M}$:
  \[
  \Expan{M}{X, b_1, \ldots, b_n} \models_L !F \text{ yen lan
    mung yen } \Struct{N} \models_L !E,
  \]
  kanthi $\Struct{N}$ iku substruktur saka $\Struct{M}$ kanthi
  !!{domain} $\Domain{N} = \Setabs{a\in \Domain{M}}{\Assign{R}{M}(a,
  b_1, \dots, b_n)}$ (waca
  \olref[bas][sub]{rem:substructure}), lan
  $\Expan{M}{X, b_1, \ldots, b_n}$ iku ekspansi $\Struct{M}$
  kang napsirake $R$, $c_1$, \dots, $c_n$ kanthi $X$, $b_1,$
  \dots, $b_n$, miturut urutane (kanthi $X \subseteq M^{n+1}$).
\end{enumerate}
\end{defn}

\begin{defn}
Yen diwenehi rong logika abstrak $\tuple{L_1, \models_{L_1}}$
lan $\tuple{L_2, \models_{L_2}}$, kita kandha yen logika kapindho
\emph{paling ora sakuwate kang kapisan anggone nyatakake},
ditulis $\tuple{L_1, \models_{L_1}} \leq
\tuple{L_2, \models_{L_2}}$, yen kanggo saben !!{language}
$\Lang{L}$ lan !!{sentence} $!E \in L_1(\Lang{L})$ ana
!!a{sentence} $!F \in L_2(\Lang{L})$ saengga
$\Mod[L](L_1){!E} = \Mod[L](L_2){!F}$. Logika
$\tuple{L_1, \models_{L_1}}$ lan
$\tuple{L_2, \models_{L_2}}$ iku \emph{ekuivalen} yen
$\tuple{L_1, \models_{L_1}} \leq \tuple{L_2,
  \models_{L_2}}$ lan $\tuple{L_2, \models_{L_2}} \leq
\tuple{L_1, \models_{L_1}}$.
\end{defn}

\begin{rem}
  Logika orde kapisan, yaiku logika abstrak $\tuple{F, \models}$,
  iku normal. Satemene, sipat-sipat ing ndhuwur umume langsung
  lumaku kanggo logika orde kapisan. Kita mung ngelingake yen
  sipat ekspansi mujudake ekstensionalitas, lan relativisasi
  !!a{sentence} $!E$ marang $\Atom{R}{x, c_1, \dots, c_n}$
  dipikolehi kanthi ngganti saben !!{subformula}
  $\lforall[x][!F]$ dadi
  $\lforall[x][(\Atom{R}{x, c_1, \dots, c_n} \lif \ !F)]$.
  Kajaba iku, yen $\tuple{L, \models_L}$ normal, mula
  $\tuple{F, \models} \leq \tuple{L, \models_{L}}$,
  kaya kang bisa dibuktekake kanthi induksi miturut
  !!{formula}s orde kapisan. Mula, tanpa ngilangake kalumrahan,
  kita bisa nganggep yen saben !!{sentence} orde kapisan
  kalebu ing saben logika normal.
\end{rem}

\end{document}

